\clearpage
{\noindent\sffamily \large
\begin{tblr}{colspec={|Q[l,m,wd=4.5cm]|Q[l,m,wd=11cm]|},
             abovesep=6pt,      % Adds 5pt spacing around rows
             belowsep=6pt       % Adds 5pt spacing around columns
	        }  \hline
    Geometry Reasoning             &{Complementary angles are two angles that add up to $90^{\circ}$.     \\[0.6cm]
                                      Angles do not have to be adjacent to be complementary.                     }                    \\ \hline
    Geometry Reasoning             &{Supplementary angles are two angles that add up to $180^{\circ}$.    \\[0.5cm]
                                      Angles do not have to be adjacent to be supplementary.                      }     	          \\ \hline
    Geometry Reasoning             & Vertically opposite angles are equal.                                                            \\ \hline
    Geometry Reasoning             & Angles on a straight line add up to $180^{\circ} $                                               \\ \hline
    Geometry Reasoning             & Angles at a point add up to $360^{\circ} $                                                       \\ \hline
    Geometry Reasoning             & Corresponding angles on parallel lines are equal.                                                \\ \hline
    Geometry Reasoning             & Alternate angles on parallel lines are equal.                                                    \\ \hline
    Geometry Reasoning             & Co-interior angles on parallel lines add to $180^{\circ} $.                                      \\ \hline
    Geometry Reasoning             & The sum of exterior angles of any polygon is $360^{\circ} $.                                     \\ \hline
\end{tblr}
}

\clearpage
{\noindent\sffamily \large
\begin{tblr}{colspec={|Q[l,m,wd=4.5cm]|Q[l,m,wd=11cm]|},
             abovesep=6pt,      % Adds 5pt spacing around rows
             belowsep=6pt       % Adds 5pt spacing around columns
	        }  \hline
    Circle Theorem                 & { \begin{tikzpicture}[scale=0.6]

    %================================================
    % Circle
    %================================================
    \draw[thick] (0,0) circle (3);

    %================================================
    % Points
    %================================================
    \coordinate (A) at (-1.55,2.57);
    \coordinate (B) at ( 1.45,2.63);
    \coordinate (C) at ( 2.50,-1.65);
    \coordinate (D) at (-2.482,-1.65);

    %================================================
    % Chords
    %================================================
    \draw[thick] (A) -- (D);
    \draw[thick] (A) -- (C);
    \draw[thick] (B) -- (D);
    \draw[thick] (B) -- (C);

    %================================================
    % Centre
    %================================================
    \fill (0,0) circle (2.2pt);

    %================================================
    % Angles x
    %================================================
    \pic[
        draw,
        thick,
        "$a$",
        angle radius=0.58cm,
        angle eccentricity=1.35
    ] {angle = D--A--C};

    \pic[
        draw,
        thick,
        "$a$",
        angle radius=0.58cm,
        angle eccentricity=1.35
    ] {angle = D--B--C};

\end{tikzpicture}
\\Angles on the same arc are equal.    }           \\ \hline
    Circle Theorem                 & The angle at the centre is double the angle at the circumference.                                \\ \hline
    Circle Theorem                 & Angle in a semicircle is $90^{\circ} $.                                                          \\ \hline
    Circle Theorem                 & Angle between tangent and radius is $90^{\circ} $.                                               \\ \hline
    Circle Theorem                 & {The alternate segment theorem.\\[15pt]
	                                  The angle between a tangent and a chord is equal to the angle in the alternate
									  segment.                                            }                                           \\ \hline
    Circle Theorem                 & Opposite angles of a cyclic quadrilateral add to $180^{\circ} $.                                 \\ \hline	
	Geometry Reasoning             & A \uwave{cyclic polygon} is a polygon whose vertices all lie on a single circle.                 \\ \hline
    Geometry Reasoning             & A cyclic quadrilateral is a quadrilateral whose vertices lie on the circumference of a circle.   \\ \hline
\end{tblr}
}

\clearpage
{\noindent\sffamily \large
\begin{tblr}{colspec={|Q[l,m,wd=4.5cm]|Q[l,m,wd=11cm]|},
             abovesep=6pt,      % Adds 5pt spacing around rows
             belowsep=6pt       % Adds 5pt spacing around columns
	        }  \hline
	Area of Parallelogram           &  $A = bh$, where $b$ is base length, $h$ is the height.                     \\ \hline
	Area of Triangle                &  $A = \frac{1}{2}bh$, where $b$ is base length, $h$ is the height.          \\ \hline
	Area of Trapezium               &  {$A = \frac{1}{2}$(Sum of lengths of parllel sides) $\times$ height\\[15pt]
	                                    $A = \frac{1}{2}(a+b)h$}.                                                  \\ \hline
	Area of Circle                  &  $A = \pi r^2$, where $r$ is the radius.                                    \\ \hline
	Circumference of Circle         &  $A = 2 \pi r$, where $r$ is the radius.                                    \\ \hline
    Midpoint of a line segment      &  $\left(\frac{x_1+x_2}{2}, \frac{y_1+y_2}{2}\right)$                        \\ \hline
    Distance between two points     &  $d=\sqrt{(x_1-x_2)^2+(y_1-y_2)^2}$                                         \\ \hline
    Gradient of a line              &  $m=\frac{y_{2}-y_1}{x_2-x_1}$                                              \\ \hline
    Parallel Lines                  &  $m_1=m_2  $                                                                \\ \hline
    Perpendicular Lines             &  $m_1 \times m_2  = -1 $                                                    \\ \hline
    Equation of a Line              &  {$y = mx +c  $   \\ $y-y_1 = m(x-x_1) $  \\ $ ax+by+c=0$ }                 \\ \hline
\end{tblr}
}
\clearpage
{\noindent\sffamily \large
\begin{tblr}{colspec={|Q[l,m,wd=4.5cm]|Q[l,m,wd=11cm]|},
             abovesep=6pt,      % Adds 5pt spacing around rows
             belowsep=6pt       % Adds 5pt spacing around columns
	        }  \hline
	Volume of Cylinder              &  $V = \pi r^2 h$, where $r$ is the radius, $h$ is the height..              \\ \hline
	Volume of Prism                 &  $V = A \times length $, where $A$ is the area of cross-section.            \\ \hline
	Volume of Pyramid               &  $V = \frac{1}{3} A h$, where $A$ is the base area, $h$ is the height.      \\ \hline
	Volume of Cone                  &  $V = \frac{1}{3} \pi r^2 h $, where $r$ is the radius, $h$ is the height.  \\ \hline
	Volume of Sphere                &  $V = \frac{4}{3} \pi r^3$, where $r$ is the radius.                        \\ \hline
	Surface Area Sphere             &  $A = 4 \pi r^2$, where $r$ is the radius.                                  \\ \hline
	Curved Surface Area of Cone     &  $A = \pi rl$, where $r$ is the radius, $l$ is the slant edge.              \\ \hline
	Curved Surface Area of Cylinder &  $A = 2 \pi r h$, where $r$ is the radius, $h$ is the height.               \\ \hline
\end{tblr}
}

\clearpage
{\noindent\sffamily \large
\begin{tblr}{colspec={|Q[l,m,wd=4.5cm]|Q[l,m,wd=11cm]|},
             abovesep=6pt,      % Adds 5pt spacing around rows
             belowsep=6pt       % Adds 5pt spacing around columns
	        }  \hline
	Sine Rule                       & {
                                       	\begin{tikzpicture}
                                       	   \draw (0,0) coordinate (A) 
                                       	               -- node[midway, below]{$c$}  (8,0) coordinate (B)
				                                       -- node[midway, right]{$a$}  (6,4) coordinate (C)
				                                       -- node[midway, above]{$b$} cycle; 
                                           \pic[draw,angle eccentricity=1.5, pic text=$A$] {angle=B--A--C};
                                           \pic[draw,angle eccentricity=1.5, pic text=$B$] {angle=C--B--A};
                                           \pic[draw,angle eccentricity=1.5, pic text=$C$] {angle=A--C--B};
                                        \end{tikzpicture}
	
	                                    $\frac{a}{\sin A}=\frac{b}{\sin B}=\frac{c}{\sin C}$\\[10pt]
	                                               $\frac{\sin A}{a}=\frac{\sin B}{b}=\frac{\sin C}{c}$  }        \\ \hline
    Cosine Rule               & {
                                       	\begin{tikzpicture}
                                       	   \draw (0,0) coordinate (C) 
                                       	               -- node[midway, below]{$b$}  (4,0) coordinate (A)
				                                       -- node[midway, right]{$c$}  (1,5) coordinate (B)
				                                       -- node[midway, left]{$a$}  cycle;
											\pic[draw,angle eccentricity=1.5, pic text=$C$] {angle=A--C--B};
                                        \end{tikzpicture}

                                     	$ c^2 = a^2 + b^2 - 2 a b \cos{C}$		}							      \\ \hline 			   
    Area of Triangle                & {
                                       	\begin{tikzpicture}
                                       	   \draw (0,0) coordinate (C) 
                                       	               -- node[midway, below]{$b$}  (8,0) coordinate (A)
				                                       --                           (6,4) coordinate (B)
				                                       -- node[midway, above]{$a$} cycle;
											\pic[draw,angle eccentricity=1.5, pic text=$C$] {angle=A--C--B};
                                        \end{tikzpicture}

                                     	$A = \frac{1}{2} a b \sin{C}$		}							          \\ \hline 			   
\end{tblr}
}

\clearpage%
\begin{center}
      \textbf{Trignometry}\\%
      \begin{tblr}{ colspec={|Q[l,wd=6cm]|Q[l,wd=6cm]|},
                    row{1} = {c},   % first row centered
				  }
					                                                                                                     \hline
                      Formulae                             & Example                             	   	              \\ \hline
                      { $ \sin{(-x)}  =-\sin{x} $ \\ 
                        $  \cos{(-x)} =\cos{x}   $ \\
                        $  \tan{(-x)} =-\tan{x}  $   }       & { $ \sin{(-30\si{\degree})} =-\sin{30\si{\degree}} $ \\
                                                                $ \cos{(-30\si{\degree})} = \cos{30\si{\degree}} $ \\
                                                      			$ \tan{(-30\si{\degree})} =-\tan{30\si{\degree}} $  }     \\ \hline
                      $\sin{(180\si{\degree}-x)}=\sin{x} $    & $\sin{150\si{\degree}} =\sin{30\si{\degree}}$             \\ \hline
      \end{tblr}
\end{center}